\section{The Formula and Its Generalisation}

In this section we prove the formula suggested by our tests, using the
results of Section~3. We give two proofs, one by induction on $m$ and one by
induction on $n$. Up to now the game always started with person $1$. In
the second part of this section we drop this assumption and find the survivor
when the game starts with any person.

\subsection{The formula \texorpdfstring{$J(n)=2\ell+1$}{J(n) = 2l + 1}}\label{sec:formula}

\subsubsection{What we observe}

In every line of Table~\ref{tab:results}, the column $J(n)$ is equal to the
column $2\ell+1$. In other words, the survivor depends only on the remainder
$\ell$.

\begin{theo}\label{thm:main}
Let $n\geq1$, written as $n=2^m+\ell$ with $0\leq\ell<2^m$. Then
\[
J(n)=2\ell+1 .
\]
\end{theo}

\subsubsection{Proof by induction on \texorpdfstring{$m$}{m}}

For $m\geq0$, let $P(m)$ be the statement
\[
P(m):\quad\text{for every integer $\ell$ with $0\leq\ell<2^m$,}\quad
J(2^m+\ell)=2\ell+1 .
\]
By Observation~\ref{obs:decomp}, every $n\geq1$ is equal to $2^m+\ell$ for
one pair $(m,\ell)$ with $0\leq\ell<2^m$. So if $P(m)$ is true for every
$m\geq0$, the theorem is proved.

\textbf{Base case: $m=0$.} Then $0\leq\ell<2^0=1$, so $\ell=0$, and
\[
J(2^0+0)=J(1)=1=2\times0+1 .
\]
So $P(0)$ is true.

\textbf{Induction step.} Let $m\geq1$, and assume that $P(m-1)$ is true. Let
$\ell$ be an integer with $0\leq\ell<2^m$. We show that
$J(2^m+\ell)=2\ell+1$, in two cases.

\emph{Case 1: $\ell$ is even}, $\ell=2k$. From $0\leq2k<2^m$ we get
$0\leq k<2^{m-1}$, and
\[
2^m+\ell=2^m+2k=2\left(2^{m-1}+k\right).
\]
By Observation~\ref{obs:even}, then by $P(m-1)$ applied to $2^{m-1}+k$,
\[
J(2^m+\ell)=2J\!\left(2^{m-1}+k\right)-1=2(2k+1)-1=4k+1=2\ell+1 .
\]

\emph{Case 2: $\ell$ is odd}, $\ell=2k+1$. From $0\leq2k+1<2^m$ we get
$2k+1\leq2^m-1$, so $k\leq2^{m-1}-1$, that is $0\leq k<2^{m-1}$. Moreover
\[
2^m+\ell=2^m+2k+1=2\left(2^{m-1}+k\right)+1 .
\]
By Observation~\ref{obs:odd}, then by $P(m-1)$ applied to $2^{m-1}+k$,
\[
J(2^m+\ell)=2J\!\left(2^{m-1}+k\right)+1=2(2k+1)+1=4k+3=2\ell+1 .
\]

In both cases $J(2^m+\ell)=2\ell+1$, so $P(m)$ is true.

\textbf{Conclusion.} $P(0)$ is true, and $P(m-1)$ implies $P(m)$ for every
$m\geq1$. By induction, $P(m)$ is true for every $m\geq0$: for every $n\geq1$,
written $n=2^m+\ell$ with $0\leq\ell<2^m$, we have $J(n)=2\ell+1$.
\hfill$\square$

\textbf{Examples.} $100=2^6+36$ gives $J(100)=73$, the value found by halving
in Section~\ref{sub:halving}. With $2026=2^{10}+1002$, the formula gives
$J(2026)=2\times1002+1=2005$ at once.

Since $\ell=n-2^m$, the formula can also be written without $\ell$:
\[
J(n)=2n-2^{m+1}+1,
\]
where $2^m$ is the largest power of $2$ not bigger than $n$. For $n=10$:
$J(10)=20-16+1=5$.

\subsubsection{A second proof, by induction on \texorpdfstring{$n$}{n}}

Relation~\eqref{eq:onestep} of Section~\ref{sec:tool}, which adds one person
at a time, gives a second proof of Theorem~\ref{thm:main}. We prove
$J(n)=2\ell+1$ by induction on $n$. For $n=1=2^0+0$ we have $J(1)=1$.

Let $n\geq2$, and assume the formula for $n-1$: if $n-1=2^{m'}+\ell'$ with
$0\leq\ell'<2^{m'}$, then $J(n-1)=2\ell'+1$. There are two cases.
\begin{itemize}
  \item If $\ell'\leq2^{m'}-2$, then $n=2^{m'}+(\ell'+1)$ with
        $\ell'+1<2^{m'}$, so $\ell=\ell'+1$. Moreover
        $J(n-1)=2\ell'+1\leq2^{m'}+\ell'-1=n-2$, so by \eqref{eq:onestep}
        $J(n)=J(n-1)+2=2\ell'+3=2\ell+1$.
  \item If $\ell'=2^{m'}-1$, then $n=2^{m'+1}$, so $\ell=0$. Here
        $J(n-1)=2^{m'+1}-1=n-1$, the last person, so by \eqref{eq:onestep}
        $J(n)=1=2\times0+1$.
\end{itemize}
In both cases $J(n)=2\ell+1$. \hfill$\square$

The first proof uses the halving of the circle ($n$ becomes about $n/2$), the
second one adds one person at a time ($n-1$ becomes $n$). They lead to the
same formula, which is a good check.

\subsubsection{What the formula explains}

The formula proves the remarks we made on Table~\ref{tab:results}.

\begin{coro}\label{cor:consequences}
Let $n\geq1$, written as $n=2^m+\ell$ with $0\leq\ell<2^m$.
\begin{enumerate}
  \item[(a)] $J(n)$ is odd, and $1\leq J(n)\leq n$.
  \item[(b)] $J(n)=1$ if and only if $n$ is a power of $2$.
  \item[(c)] $J(n)=n$ if and only if $n=2^{m+1}-1$, as for $n=1,3,7,15$.
  \item[(d)] $J(n+1)=J(n)+2$ if $n+1$ is not a power of $2$, and $J(n+1)=1$
        if $n+1$ is a power of $2$.
\end{enumerate}
\end{coro}

\begin{proof}
By Theorem~\ref{thm:main}, $J(n)=2\ell+1$.
(a) The number $2\ell+1$ is odd and at least $1$. Moreover
$2\ell+1\leq2^m+\ell$ is the same as $\ell\leq2^m-1$, which is true.
(b) $2\ell+1=1$ if and only if $\ell=0$, that is $n=2^m$.
(c) $2\ell+1=2^m+\ell$ if and only if $\ell=2^m-1$, that is $n=2^{m+1}-1$.
(d) If $n+1$ is not a power of $2$, then $\ell+1<2^m$ (otherwise
$n+1=2^{m+1}$), so $n+1=2^m+(\ell+1)$ with the same $m$, and
$J(n+1)=2(\ell+1)+1=J(n)+2$. If $n+1$ is a power of $2$, then $J(n+1)=1$ by
(b).
\end{proof}

Point (a) agrees with Step~1 of Section~\ref{sec:even}: every even number
leaves during the first round. Point (d) is relation~\eqref{eq:onestep} again.
The table below summarises these points. On each block $2^m\leq n<2^{m+1}$,
the survivor starts again at $1$ and goes through the odd numbers in
increasing order, up to the last person $2^{m+1}-1$.
\begin{center}
\renewcommand{\arraystretch}{1.2}
\begin{tabular}{|c|c|c|}
\hline
$m$ & block of $n$ & values of $J(n)$, in order\\ \hline
$0$ & $1$ & $1$\\
$1$ & $2,3$ & $1,3$\\
$2$ & $4,\dots,7$ & $1,3,5,7$\\
$3$ & $8,\dots,15$ & $1,3,5,\dots,15$\\
$4$ & $16,\dots,31$ & $1,3,5,\dots,31$\\ \hline
$m$ & $2^m,\dots,2^{m+1}-1$ & $1,3,5,\dots,2^{m+1}-1$\\ \hline
\end{tabular}
\end{center}
A block has $2^m$ values of $n$, and there are exactly $2^m$ odd numbers
between $1$ and $2^{m+1}-1$, so each odd number appears once in each block.

\subsection{When the game starts with person \texorpdfstring{$s$}{s}}\label{sec:general}

So far, person $1$ was always the first person skipped
(Section~\ref{sec:assumption}). In a real game we may not know who starts. So
we ask: \emph{if the game starts with person $s$, who survives?}

\begin{defi}
Let $n\geq1$ and $1\leq s\leq n$. We write $J_s(n)$ for the survivor of the
game on the circle $1,2,\dots,n$ that starts with person $s$: we skip $s$,
eliminate the next person, and go on with the rules of
Definition~\ref{def:game}. In particular $J_1(n)=J(n)$.
\end{defi}

\subsubsection{\texorpdfstring{An example: $n=10$ and $s=4$}{An example: n = 10 and s = 4}}

The rounds are now counted from person $4$.
\begin{itemize}
  \item Round 1 skips $4$ and eliminates $5,7,9,1,3$. The game is back at
        person $4$.
  \item Round 2 skips $4$ and eliminates $6$, skips $8$ and eliminates $10$,
        skips $2$ and eliminates $4$.
  \item Only $8$ and $2$ are left. Round 3 skips $8$ and eliminates $2$.
\end{itemize}
The survivor is $8$, so $J_4(10)=8$ (Figure~\ref{fig:s4}).

\begin{figure}[!ht]
\centering
\setlength{\tabcolsep}{4pt}
\begin{tabular}{cccc}
\begin{tikzpicture}[baseline=(current bounding box.center)]
  \foreach \p in {1,...,10}
    \node[person] (q\p) at ({90-36*(\p-1)}:1.05) {\p};
  \draw[dashed,thick] (q4) circle ({0.5*\nodesize+0.8mm});
\end{tikzpicture}
& \gamecircle{10}{1.05}{}{5,7,9,1,3}{0}
& \gamecircle{10}{1.05}{5,7,9,1,3}{6,10,4}{0}
& \gamecircle{10}{1.05}{5,7,9,1,3,6,10,4}{2}{8}\\[3pt]
\footnotesize start: $s=4$ & \footnotesize Round 1: 5, 7, 9, 1, 3
& \footnotesize Round 2: 6, 10, 4 & \footnotesize Round 3: 2
\end{tabular}
\caption{The game with $10$ people that starts with person $4$ (dashed):
$J_4(10)=8$.}
\label{fig:s4}
\end{figure}

When person $1$ starts, the survivor is $J(10)=5$. Here the game starts three
places later, and the survivor is also three places later: $5+3=8$.

\subsubsection{Counting places from the first person}

The rule \emph{skip one, eliminate the next} only uses the order of the people
around the circle. So we number the \emph{places} from the person who starts:
place $1$ is person $s$, place $2$ is the person after him, and so on. For
$n=10$ and $s=4$:
\begin{center}
\renewcommand{\arraystretch}{1.2}
\begin{tabular}{|c|c|c|c|c|c|c|c|c|c|c|}
\hline
place & 1&2&3&4&5&6&7&8&9&10\\ \hline
person & 4&5&6&7&8&9&10&1&2&3\\ \hline
\end{tabular}
\end{center}
By Lemma~\ref{lem:places}, the survivor is at place $J(10)=5$, that is
person $8$. In general the survivor is at place $J(n)$, and we only need to
know who stands there. Going forward from $s$ we meet $s+1$, $s+2$, \dots, up
to person $n$, and then we come back to person $1$. So the person at place $i$
is $s+i-1$ if $s+i-1\leq n$, and $s+i-1-n$ otherwise. For instance, place $8$
gives $4+8-1=11>10$, so person $11-10=1$.

\subsubsection{The general formula}

Since the survivor is at place $J(n)=2\ell+1$, we obtain the following
result.

\begin{theo}\label{thm:general}
Let $n\geq1$ and $1\leq s\leq n$, and write $n=2^m+\ell$ with $0\leq\ell<2^m$.
If the game starts with person $s$, the survivor is
\[
J_s(n)=
\begin{cases}
s+2\ell & \text{if } s+2\ell\leq n,\\
s+2\ell-n & \text{if } s+2\ell>n.
\end{cases}
\]
In words: the survivor is the person $2\ell$ places after $s$ around the
circle.
\end{theo}

\begin{proof}
The game that starts with $s$ starts from the list
$(a_1,\dots,a_n)=(s,s+1,\dots,n,1,\dots,s-1)$ of the people in the order of
the circle, where $a_i$ is the person at place $i$. By
Lemma~\ref{lem:places}, the survivor is $a_{J(n)}$, the person at place
$J(n)$, and $J(n)=2\ell+1$ by Theorem~\ref{thm:main}. The person at
place $2\ell+1$ is $s+2\ell$ if $s+2\ell\leq n$, and $s+2\ell-n$ otherwise.

It remains to check that one subtraction of $n$ is always enough, that is
$s+2\ell-n\leq n$. Since $\ell<2^m$, we have $2\ell<2^m+\ell=n$, and with
$s\leq n$ we get $s+2\ell<2n$.
\end{proof}

\textbf{Another way to read the formula.} Since $J(n)=2\ell+1$, we have
$s+2\ell=J(n)+(s-1)$. So
\[
J_s(n)=J(n)+(s-1),\quad\text{minus $n$ if this is bigger than $n$.}
\]
Starting the game $s-1$ places later moves the survivor $s-1$ places later.
This is what we saw for $n=10$ and $s=4$: $J_4(10)=5+3=8$.

\textbf{Checks.} For $s=1$ the formula gives $J_1(n)=1+2\ell=J(n)$, as it should. For
$n=10=2^3+2$ we have $\ell=2$, so the survivor is $4$ places after $s$:
\begin{center}
\renewcommand{\arraystretch}{1.2}
\begin{tabular}{|c|c|c|c|c|c|c|c|c|c|c|}
\hline
$s$ & 1&2&3&4&5&6&7&8&9&10\\ \hline
$J_s(10)$ & 5&6&7&8&9&10&1&2&3&4\\ \hline
\end{tabular}
\end{center}
From $s=7$ on, $s+4>10$ and we subtract $10$: for instance $J_8(10)=8+4-10=2$.
We also checked the formula with a computer simulation of the game, for every
$n\leq200$ and every starting person $s$.
